\section{Every fixed height}\label{sec:height}

Fix the height $h$ and read the grid from left to right. Its blackout is
the \emph{column word} $w=B_1\cdots B_m$, whose letter $B_i$ records the
black rows in column $i$. Write $|w|=m$ and $S(w)$ for its blackout, and
let $\mathcal L_h$ be the valid words, including the empty word. We first
show that validity needs only finite memory. A bounded counter will then
pick out the words with the most black points.

\begin{theorem}[regularity and rational series]\label{thm:regular}
For every $h\ge1$ the language $\mathcal L_h$ is regular, and $G_h(z,u)=\sum_{w\in\mathcal L_h}z^{|w|}u^{|S(w)|}$ is rational in $z$ and $u$.
\end{theorem}

\begin{proof}
The \emph{trace} of a rectangle records, in each column, the rows of its
corners. It determines the corner set. By Lemma~\ref{lem:diff}, two distinct
traces $Y,Y'$ give the same presentation exactly when
$Y_i\triangle Y'_i\subseteq B_i$ in every column.

The possible traces form a regular language. For an axis-aligned rectangle,
choose $r<s$, read a column $\{r,s\}$, then another such column after any
number of empty columns. There can also be empty columns before and after.
For a tilted rectangle, Lemmas~\ref{lem:leftmost} and~\ref{lem:span} leave
only finitely many templates~\eqref{eq:tilt-param}, each of span at most
$W_h$, again with arbitrary empty columns on either side.

An automaton can therefore read $w$ while guessing two rectangle traces.
It checks the containment above and remembers whether the traces have
differed anywhere. Accepting when they differ recognizes the invalid words.
Regular languages are closed under complementation, so $\mathcal L_h$ is
regular too.

Choose a deterministic automaton for $\mathcal L_h$ and give a transition
reading $B$ the weight $u^{|B|}$. If $M(u)$ is its weighted transition
matrix, $\alpha$ its start vector and $\beta$ its accepting vector, then
\[
 G_h(z,u)=\alpha(I-zM(u))^{-1}\beta .
\]
Each word is counted once because the automaton is deterministic. Cramer's
rule makes the expression rational~\cite[Proposition~V.6]{flajolet2009}.
\end{proof}

\begin{theorem}[fixed height]\label{thm:periodic}
Let $h\ge2$. (a) There are integers $m_0\ge3$, $p\ge1$, and $k_0,\dots,k_{p-1}$ such that $f(h,m)=m+k_{m\bmod p}$ for every $m\ge m_0$. (b) The series $\sum_{m\ge0}a_h(m)z^m$ is rational, with the conventions of Section~\ref{sec:conv}, and the words of optimum blackouts form a regular language.
\end{theorem}

\begin{proof}
Put $E(w)=|S(w)|-|w|$, the number of black points beyond one per column.
Every consecutive block $v$ of a valid word is valid on its own grid.
Theorem~\ref{thm:rowbound} bounds $E(v)$ by $q_h$ when $|v|\ge3$.
For a shorter block, $|S(v)|\le h|v|$ gives the common bound
\[
 E(v)\le C_h:=\max(q_h,2h-2).
\]
If a valid word $w=uv$ has excess $k$, applying this to the prefix and
suffix gives
\[
 k-C_h\le E(u)=k-E(v)\le C_h.
\]
Thus, for a fixed $k$, only finitely many prefix excesses need to be
remembered. Add this counter to an automaton for $\mathcal L_h$, and
accept when the final excess is $k$. It recognizes the regular language
$\mathcal L_{h,k}$ of valid words with that excess.

The set of lengths of a regular language is eventually periodic. For
$m\ge3$, Lemma~\ref{lem:cross} and Theorem~\ref{thm:rowbound} put the
optimal excess between $h-1$ and $q_h$. There are only finitely many such
$k$, so choose one common period $p$ and threshold $m_0$ for their length
sets. The largest available excess then depends only on $m\bmod p$.
Calling it $k_r$ proves (a).

For each residue $r$, keep the words of $\mathcal L_{h,k_r}$ whose lengths
are at least $m_0$ and congruent to $r$ modulo $p$. This is still a regular
language. Their union, together with the finitely many shorter optimum
words, is exactly the language of optimum blackouts. Its counting series
is rational by the same transfer-matrix argument as above, proving (b).
\end{proof}
